% ==================================================
% Fundamental Theorem of Finite Calculus
% Discrete Calculus
% ==================================================

\section*{Fundamental Theorem of Finite Calculus}

\begin{tcolorbox}[colback=thmbox, colframe=thmborder, arc=2pt,
  left=6pt, right=6pt, top=4pt, bottom=4pt,
  title={\small\textbf{Theorem (Fundamental Theorem of Finite Calculus)}},
  fonttitle=\small\bfseries]
Let \(F\) be a discrete antiderivative of \(f\), so that
\(\Delta F(n) = f(n)\) for all \(n\). Then
\[
\sum_{k=a}^{b-1} f(k) = F(b) - F(a).
\]
\end{tcolorbox}
\label{thm:ftfc}

\begin{remark*}[Fully quantified form]
\(\forall f,F : \mathbb{Z}\to\mathbb{R},\;
\forall a,b \in \mathbb{Z},\; a \le b:\;
(\forall n,\; \Delta F(n) = f(n))
\Rightarrow
\displaystyle\sum_{k=a}^{b-1} f(k) = F(b) - F(a)\).
\end{remark*}

\begin{proof}
Since \(f(k) = \Delta F(k)\) for all \(k\), substituting gives
\[
\sum_{k=a}^{b-1} f(k)
=
\sum_{k=a}^{b-1} \Delta F(k).
\]
By the Telescoping Identity (Theorem~\ref{thm:telescoping}),
\[
\sum_{k=a}^{b-1} \Delta F(k) = F(b) - F(a). \qedhere
\]
\end{proof}

\begin{remark*}[Analogy with continuous calculus]
This theorem is the discrete analogue of the Fundamental Theorem of
Calculus,
\[
\int_a^b f'(x)\,dx = f(b)-f(a).
\]
In continuous calculus, differentiation and integration are inverse
operations. In discrete calculus, the analogous inverse pair is the
difference operator \(\Delta\) and summation. Any sum can be evaluated
by finding a discrete antiderivative of the summand — just as any
definite integral can be evaluated by finding an antiderivative.
\end{remark*}

\begin{remark*}[Practical consequence]
The theorem reduces the problem of computing \(\sum_{k=a}^{b-1} f(k)\)
to the algebraic problem of finding \(F\) with \(\Delta F = f\). Combined
with the falling factorial rules developed in the next section, this
gives a systematic method for evaluating a wide class of sums in
closed form.
\end{remark*}
